\documentclass[11pt]{article}
\usepackage{amsmath,amssymb,amsthm,hyperref}
\newtheorem{theorem}{Theorem}
\newcommand{\E}{\mathbb E}
\emergencystretch=2em
\title{Local scalar-transfer rigidity at strictly constrained endpoints}
\author{Research note; independently reviewed}
\date{30 September 2026}
\begin{document}\maketitle
This extends the regular interior-boundary result in
\path{../../independent_directions/scalar_transfer_boundary.tex}.
The statement explicitly assumes strict endpoint signs. It makes no
claim about endpoint strata where the normal derivative vanishes.

\begin{theorem}
Fix $m\ge2$ and a scalar array $a:[K]\to[0,1]$, nondecreasing,
with at most $m-1$ distinct levels in $(0,1)$. Give its rows fixed
positive probability weights. Endpoint levels $0$ and $1$ may also
occur. Suppose a degree-$m$ polynomial
\[
 P(t)=\sigma t^m+\sum_{s=0}^{m-1}\gamma_s t^s,
 \qquad \sigma\in\{-1,1\},
\]
has the following properties:
\begin{itemize}
\item $P'(a_l)=0$ at every interior level;
\item $P''(a_l)>0$ at every interior level containing at least two rows;
\item $P'(0)>0$ if the level $0$ occurs, and $P'(1)<0$ if level $1$ occurs.
\end{itemize}
Then, in a neighborhood of the scalar array with all columns equal
to $a$, every nondecreasing $[0,1]$-valued array $x_1,\ldots,x_m$
with exchangeable mixed moments
\[
 \E\prod_{j\in S}x_j=\mu_{|S|}
\]
and with $\mu_s=\E a^s$ for $1\le s<m$ satisfies
\[
 \sigma(\mu_m-\E a^m)\ge0.
\]
In the equality case all endpoint blocks are unchanged; on every
interior block all columns have mean $a_l$, and at most one column is
nonconstant. Thus all mixed moments through degree $m$ equal those
of the original scalar array.
\end{theorem}
\begin{proof}
All constants depend only on the fixed scalar array and $m$.
Let $\delta=\max_{j,q}|x_j(q)-a(q)|$.
On each interior level block write
\[
 x_j=a_l+b_{jl}+e_j,\qquad \E_l e_j=0,
\]
where $\E_l$ is the conditional weighted mean on that block. Put
\[
 C_{lij}=\E_l(e_i e_j)\ge0,
 \qquad S=\sum_{l\text{ interior}}\sum_{i<j}C_{lij},
 \qquad B=\max_{j,l\text{ interior}}|b_{jl}|.
\]
If there are no interior levels, take $B=S=0$.
For endpoint blocks put $b_{j0}=\E_0 x_j\ge0$ and
$b_{j1}=\E_1(x_j-1)\le0$, omitting absent blocks, and set
\[
 T=\sum_j(b_{j0}-b_{j1})\ge0.
\]

We use the elementary inequality
$\E_l|fg|\le3\E_l(fg)$ for two nondecreasing mean-zero functions.
Indeed, represent them as positive combinations of centered steps
$h_s(u)=\mathbf1_{u\ge s}-(1-s)$ after embedding the weighted block
in $[0,1]$. For $s\le t$,
\[
 \E h_s h_t=s(1-t),\qquad
 \E|h_s h_t|=s(1-t)[1+2(t-s)]\le3s(1-t).
\]
Consequently, any product of at least two distinct interior errors
has absolute mean at most $3(2\delta)^{s-2}C_{lij}$, where $s$ is
its number of factors and $i,j$ are any two of them.

The lower mixed moments of the block-constant interior means differ
from their prescribed scalar values by $O(S+T)$. Here interior
fluctuations contribute $O(S)$, while endpoint blocks contribute
$O(T)$ because their deviations have one fixed sign and are bounded
by $\delta$.
The derivative of the lower-moment map in the interior variables
$b_{jl}$ is injective. A kernel vector would satisfy, for every
$s$-element column set $I$, $1\le s<m$,
\[
 \sum_{j\in I}\sum_{l\text{ interior}} w_l a_l^{s-1}b_{jl}=0,
\]
where $w_l$ is the fixed block weight. Comparing sets differing in one
index makes each inner sum zero. The first $r$ equations, where $r$
is the number of interior levels, form an invertible weighted
Vandermonde system for each column. A bounded left inverse and Taylor
expansion therefore give
$B\le C(S+T+B^2)$, hence
\[
 B=O(S+T)
\]
after shrinking the neighborhood.

The symmetric multiaffine polarization of $P$ at a row
$(a_l+d_1,\ldots,a_l+d_m)$ is
\[
 P(a_l)+\sum_{s=1}^m\frac{P^{(s)}(a_l)}{s!}
                  \frac{e_s(d_1,\ldots,d_m)}{\binom ms}.
\]
On interior blocks the linear term vanishes. Their quadratic main
term is
\[
 \sum_{l\text{ interior}} w_l\frac{P''(a_l)}{m(m-1)}
                    \sum_{i<j}C_{lij}.
\]
Singleton blocks contribute no covariance. Thus this term is at least
$cS$ if $S$ can be nonzero; all errors, including products of block
mean shifts, are $O(\delta S+(S+T)^2)$.

On an endpoint block, every deviation $d_j=x_j-a_l$ has the same
sign. A product of $s\ge2$ of them has absolute mean at most
$\delta^{s-1}\E_l|d_j|$ for any chosen factor. Hence the endpoint
contribution is its linear term plus $O(\delta T)$. By the strict
endpoint signs and the positive block weights, its linear term is
at least $cT$.
Combining the two types of blocks gives
\[
 \E\operatorname{Pol}(P)(x_1,\ldots,x_m)-\E P(a)
 \ge c(S+T)-C\bigl(\delta(S+T)+(S+T)^2\bigr)\ge0
\]
for sufficiently small $\delta$, with strict inequality if $S+T>0$.
Exchangeability and the prescribed lower moments identify the left
side as $\sigma(\mu_m-\E a^m)$.

Equality forces $S=T=0$, hence $B=0$. The endpoint blocks are fixed
because their mean one-sided deviations vanish. On an interior block,
zero covariance between two monotone centered functions forces one to
be constant: the monotone covariance formula has a strictly positive
first--last term if both vary. Thus at most one column varies on that
block, and its mean is the scalar level. Every distinct-coordinate
product has the scalar block mean, proving the equality statement.
\end{proof}

At a scalar constrained local minimum, the usual one-sided endpoint
conditions give weak signs $P'(0)\ge0$ and $P'(1)\le0$. The theorem
covers the strictly complementary case. Vanishing endpoint derivatives,
singular interior strata, and global components remain outside its
scope. In particular this note does not establish general scalar
transfer or a universal quadratic per-die lower bound.
\end{document}
